\documentclass[10pt,a4paper]{amsart}
\usepackage[margin=19mm]{geometry}
\usepackage{amsmath,amssymb,mathtools}
\usepackage[scr=rsfs,cal=euler]{mathalfa}
\usepackage{xfrac}
\usepackage[unicode,hidelinks,bookmarks=false]{hyperref}
\newcommand{\ie}{i.e.\ }
\newcommand{\cf}{cf.\ }
\newcommand{\resp}{respectively\ }
\newcommand{\Subsection}{\subsection}
\providecommand{\nobreakdash}{\nobreak\hbox{-}}
\setlength{\emergencystretch}{2em}
\multlinegap0pt
\usepackage{tikz}
\usetikzlibrary{cd}
\usepackage{mathtools}
\usepackage{amssymb}
\usepackage{xcolor}
\usepackage{enumerate}
\usepackage{tikz}
\newtheorem{theorem}{Theorem}
\DeclareMathOperator{\THH}{THH}
\newcommand{\id}{\text{id}}
\newcommand{\smashy}{\wedge}
\newcommand{\wedgey}{\vee}
\title{The algebraic structure of twisted topological Hochschild homology}
\author{Danika Van Niel}
\date{Source: arXiv:2602.02423v1 (CC BY 4.0)}
\begin{document}
\maketitle

\noindent\emph{Source and licence.} The algebraic structure of twisted topological Hochschild homology. Version arXiv:2602.02423v1 (CC BY 4.0), distributed by arXiv under the Creative Commons Attribution 4.0 International licence, http://creativecommons.org/licenses/by/4.0/, as recorded in the arXiv licence metadata for this exact version and shown on the abstract page https://arxiv.org/abs/2602.02423v1. Full licence notice: \texttt{licenses/arxiv-2602.02423-CC-BY-4.0.txt}.\par\smallskip

\noindent\textit{Editorial scope.} Counted unit: the article's theorem THHAlgebra (source0.tex lines 1871--1873). The article's complete original proof of it is delivered with it (source0.tex lines 1875--1924, one \texttt{proof} environment of 50 lines). No mathematics is written, repaired or re-derived: the statement and the proof are the article's own lines, in the article's own notation, and no retained line is substituted or re-flowed.  No further source-local context is needed: every cross-reference of every retained line resolves inside the two delivered spans.  Nothing was omitted from any retained span, so the package carries no omission record.  No retained line contains a citation, so the wrapper carries no bibliography and no bibliography entry.  The retained lines contain the article's own diagram code, which the wrapper typesets with the standard tikz packages; no image file is delivered and the package declares no asset dependency.  Editorial formatting only, in addition: an AMS \texttt{amsart} preamble that restates the article's own declarations and macros used by the retained lines, namely the environments \texttt{theorem} on separate counters, together with the standard packages those lines need.  The retained environments print in this document as Theorem 1 in order of appearance, whereas the article prints the same items with the numbering of its own full document.  The complete native source of the article as distributed in the arXiv e-print for this exact version is bundled unchanged as \texttt{sources/193852/source0.tex}.
\begin{theorem} \label{THHAlgebra}
Let $p$ be prime. For a commutative ring $C_p$-spectrum $R$, $\THH_{C_p}(R)$ is a commutative $R$-algebra in the category of commutative ring $C_p$-spectra.
\end{theorem}

\begin{proof} We begin by checking associativity of the product map $\phi\colon \THH_{C_p}(R) \smashy_R \THH_{C_p}(R) \to \THH_{C_p}(R)$. For ease of notation, let $T \coloneqq \THH_{C_p}(R)$. We need to verify that the following diagram commutes:
% https://q.uiver.app/#q=WzAsNCxbMCwwLCJUIFxcc21hc2h5X1IgVCBcXHNtYXNoeV9SIFQiXSxbMiwwLCJUIFxcc21hc2h5X1IgVCJdLFswLDIsIlQgXFxzbWFzaHlfUiBUIl0sWzIsMiwiVCJdLFswLDEsIlxcaWQgXFxzbWFzaHkgXFxwaGkiXSxbMCwyLCJcXHBoaSBcXHNtYXNoeSBcXGlkIiwyXSxbMiwzLCJcXHBoaSIsMl0sWzEsMywiXFxwaGkiXV0=
\[\begin{tikzcd}[column sep=small,row sep=small]
	{T \smashy_R T \smashy_R T} && {T \smashy_R T} \\
	\\
	{T \smashy_R T} && T.
	\arrow["{\id \smashy \phi}", from=1-1, to=1-3]
	\arrow["{\phi \smashy \id}"', from=1-1, to=3-1]
	\arrow["\phi", from=1-3, to=3-3]
	\arrow["\phi"', from=3-1, to=3-3]
\end{tikzcd}\]
\noindent It is sufficient to show that the following diagram of $C_p$-simplicial spaces commutes:
% https://q.uiver.app/#q=WzAsNCxbMCwwLCJwU14xX1xcYnVsbGV0IFxcd2VkZ2V5X3tDX3B9IHBTXjFfXFxidWxsZXQgXFx3ZWRnZXlfe0NfcH0gcFNeMV9cXGJ1bGxldCJdLFsyLDAsInBTXjFfXFxidWxsZXQgXFx3ZWRnZXlfe0NfcH0gcFNeMV9cXGJ1bGxldCJdLFswLDIsInBTXjFfXFxidWxsZXQgXFx3ZWRnZXlfe0NfcH0gcFNeMV9cXGJ1bGxldCJdLFsyLDIsInBTXjFfXFxidWxsZXQiXSxbMCwxLCJcXGlkIFxcd2VkZ2V5IFxccGhpIl0sWzAsMiwiXFxwaGkgXFx3ZWRnZXkgXFxpZCIsMl0sWzIsMywiXFxwaGkiLDJdLFsxLDMsIlxccGhpIl1d
\[\begin{tikzcd}[column sep=small,row sep=small]
	{pS^1_\bullet \wedgey_{C_p} pS^1_\bullet \wedgey_{C_p} pS^1_\bullet} && {pS^1_\bullet \wedgey_{C_p} pS^1_\bullet} \\
	\\
	{pS^1_\bullet \wedgey_{C_p} pS^1_\bullet} && {pS^1_\bullet.}
	\arrow["{\id \wedgey \phi}", from=1-1, to=1-3]
	\arrow["{\phi \wedgey \id}"', from=1-1, to=3-1]
	\arrow["\phi", from=1-3, to=3-3]
	\arrow["\phi"', from=3-1, to=3-3]
\end{tikzcd}\]
\noindent The maps $\id \wedgey \phi$ and $\phi \wedgey \id$ fold the inner two and outer two copies of $pS^1_\bullet$ together, respectively, where $\phi$ folds the remaining two copies of $pS^1_\bullet$ together. Therefore this diagram commutes.

To check unitality and commutativity of the product map, we need to show that the following diagrams commute:
% https://q.uiver.app/#q=WzAsNyxbMCwwLCJwU14xX1xcYnVsbGV0IFxcd2VkZ2V5X3tDX3B9IENfcCJdLFsyLDAsInBTXjFfXFxidWxsZXQgXFx3ZWRnZXlfe0NfcH0gcFNeMV9cXGJ1bGxldCJdLFs0LDAsIkNfcCBcXHdlZGdleV97Q19wfSBwU14xX1xcYnVsbGV0Il0sWzIsMywicFNeMV9cXGJ1bGxldCJdLFswLDUsInBTXjFfXFxidWxsZXQgXFx3ZWRnZXlfe0NfcH0gcFNeMV9cXGJ1bGxldCJdLFs0LDUsInBTXjFfXFxidWxsZXQgXFx3ZWRnZXlfe0NfcH0gcFNeMV9cXGJ1bGxldCJdLFsyLDgsInBTXjFfXFxidWxsZXQiXSxbMCwxLCJcXGlkIFxcd2VkZ2V5IFxcZXRhIl0sWzIsMSwiXFxldGEgXFx3ZWRnZXkgXFxpZCIsMl0sWzAsMywiXFxzaW1lcSIsMl0sWzIsMywiXFxzaW1lcSJdLFsxLDMsIlxccGhpIl0sWzUsNiwiXFxwaGkiXSxbNCw2LCJcXHBoaSIsMl0sWzQsNSwiXFx0YXUiXV0=
\[\begin{tikzcd}[column sep=small,row sep=tiny]
	{pS^1_\bullet \wedgey_{C_p} C_p} && {pS^1_\bullet \wedgey_{C_p} pS^1_\bullet} && {C_p \wedgey_{C_p} pS^1_\bullet} \\
	\\
	\\
	&& {pS^1_\bullet} \\
	{pS^1_\bullet \wedgey_{C_p} pS^1_\bullet} &&&& {pS^1_\bullet \wedgey_{C_p} pS^1_\bullet} \\
	\\
	\\
	&& {pS^1_\bullet}
	\arrow["{\id \wedgey \eta}", from=1-1, to=1-3]
	\arrow["\simeq"', from=1-1, to=4-3]
	\arrow["\phi", from=1-3, to=4-3]
	\arrow["{\eta \wedgey \id}"', from=1-5, to=1-3]
	\arrow["\simeq", from=1-5, to=4-3]
	\arrow["\tau", from=5-1, to=5-5]
	\arrow["\phi"', from=5-1, to=8-3]
	\arrow["\phi", from=5-5, to=8-3]
\end{tikzcd}\]
\noindent where $pS^1_\bullet \wedgey_{C_p} C_p \cong pS^1_\bullet$ is the pushout of the span $pS^1_\bullet \leftarrow C_p \to C_p$. The map $\id \wedgey \eta\colon pS^1_\bullet \wedgey_{C_p} C_p \to pS^1_\bullet \wedgey_{C_p} pS^1_\bullet$ is the identity on $pS^1_\bullet$ and includes $C_p$ into the second copy of $pS^1_\bullet$, similarly $\eta \wedgey \id\colon C_p \wedgey_{C_p} pS^1_\bullet \to pS^1_\bullet \wedgey_{C_p} pS^1_\bullet$ is the identity on $pS^1_\bullet$ and includes $C_p$ into the first copy of $pS^1_\bullet$. The map $\tau$ swaps the first and second copies of $pS^1_\bullet$. 

For the unitality diagram, note that $\phi\colon pS^1_\bullet \wedgey_{C_p} pS^1_\bullet \to pS^1_\bullet$ is the fold map and the maps $\id \wedgey \eta$ and $\eta \wedgey \id$ both have an image of one copy of $pS^1_\bullet$, so this diagram commutes. For the commutativity diagram, the fold map has the same image no matter the position of the two copies of $pS^1_\bullet$.

Therefore, $\THH_{C_p}(R)$ is a commutative $R$-algebra.
\end{proof}

\end{document}
